\begin{exo}{}
	$ABCD$ est un carré dont le côté mesure $4$ cm. $DFC$ et $BCH$ sont des triangles équilatéraux.

	\begin{center}
		\begin{tikzpicture}[scale=0.5, >=Stealth]
			\coordinate (A) at (0,0);
			\coordinate (B) at (4,0);
			\coordinate (C) at (4,4);
			\coordinate (D) at (0,4);
			\coordinate (F) at (2, {4 + 2*sqrt(3)});
			\coordinate (H) at ({4 - 2*sqrt(3)}, 2);
			\draw[thick] (A) -- (B) -- (C) -- (D) -- cycle;
			\draw[thick] (D) -- (F) -- (C);
			\draw[thick] (B) -- (H) -- (C);
			\node[below left] at (A) {$A$};
			\node[below right] at (B) {$B$};
			\node[right] at (C) {$C$};
			\node[left] at (D) {$D$};
			\node[above] at (F) {$F$};
			\node[above] at (H) {$H$};
			\foreach \pt in {A, B, C, D, F, H} { \filldraw (\pt) circle (3pt); }
		\end{tikzpicture}
	\end{center}

	\begin{enumerate}
		\item Calculer la hauteur de chaque triangle équilatéral.
		\item Démontrer que les points $A$, $H$ et $F$ sont alignés.
	\end{enumerate}

	\textit{Aide : On pourra se placer dans un repère convenablement choisi.}
\end{exo}
